\documentclass[12pt, a4paper, oneside]{report}

% ─── Packages ───────────────────────────────────────────────
\usepackage[top=1in, bottom=1in, left=1.25in, right=1in]{geometry}
\usepackage{times}
\usepackage{setspace}
\usepackage{titlesec}
\usepackage{fancyhdr}
\usepackage{graphicx}
\usepackage{hyperref}
\usepackage{tocloft}
\usepackage{parskip}
\usepackage{enumitem}
\usepackage{booktabs}
\usepackage{array}
\usepackage{tabularx}

\hypersetup{
    colorlinks=true,
    linkcolor=black,
    urlcolor=blue,
    citecolor=black
}

\onehalfspacing

% ─── Header/Footer ──────────────────────────────────────────
\pagestyle{fancy}
\fancyhf{}
\fancyfoot[L]{\small NMIET, Department of Computer Engineering 2025-26}
\fancyfoot[R]{\small \thepage}
\renewcommand{\headrulewidth}{0pt}
\renewcommand{\footrulewidth}{0.4pt}

\fancypagestyle{plain}{
    \fancyhf{}
    \fancyfoot[L]{\small NMIET, Department of Computer Engineering 2025-26}
    \fancyfoot[R]{\small \thepage}
    \renewcommand{\headrulewidth}{0pt}
    \renewcommand{\footrulewidth}{0.4pt}
}

% ─── Chapter Title Format ───────────────────────────────────
\titleformat{\chapter}[display]
  {\normalfont\Large\bfseries\centering}
  {CHAPTER \Roman{chapter}}{10pt}{\Large\centering}
\titlespacing*{\chapter}{0pt}{20pt}{15pt}

\titleformat{\section}{\normalfont\normalsize\bfseries}{\thesection}{1em}{}
\titleformat{\subsection}{\normalfont\normalsize\bfseries}{\thesubsection}{1em}{}

% ════════════════════════════════════════════════════════════
\begin{document}

% ════════════════════════════════════════════════════════════
%  PAGE 1 — TITLE PAGE
% ════════════════════════════════════════════════════════════
\begin{titlepage}
\thispagestyle{empty}
\centering

\vspace*{0.5cm}
{\large An}\\[0.3cm]
{\LARGE \textit{Internship Report}}\\[0.5cm]

{\normalsize SUBMITTED TO THE SAVITRIBAI PHULE PUNE UNIVERSITY, PUNE}\\
{\normalsize IN THE PARTIAL FULFILLMENT OF THE REQUIREMENTS}\\
{\normalsize FOR THE AWARD OF THE DEGREE}\\[0.4cm]

{\normalsize \textit{of}}\\[0.3cm]

{\large \textbf{BACHELOR OF ENGINEERING}}\\[0.3cm]
{\normalsize \textit{in}}\\[0.3cm]
{\large \textbf{COMPUTER ENGINEERING}}\\[0.8cm]

{\normalsize SUBMITTED BY}\\
{\normalsize \textbf{ADITYA SHAILESH BHARATE}}\\[0.8cm]

{\normalsize \textbf{UNDER THE GUIDANCE OF}}\\
{\normalsize PROF -- ASHWINI UTTURE}\\[0.8cm]

{\normalsize DEPARTMENT OF COMPUTER ENGINEERING}\\
{\normalsize (NBA ACCREDITED)}\\[0.8cm]

% ── Replace 'nmiet_logo.png' with your actual logo file ──
\includegraphics[width=0.22\textwidth]{nmiet_logo.png}\\[0.8cm]

{\normalsize NUTAN MAHARASHTRA INSTITUTE OF ENGINEERING \& TECHNOLOGY}\\
{\normalsize \textbf{TALEGAON DABHADE TAL.MAVAL, PUNE, 410507}}\\
{\normalsize SAVITRIBAIPHULEPUNEUNIVERSITY}\\
{\normalsize 2025 --2026}

\end{titlepage}


% ════════════════════════════════════════════════════════════
%  PAGE 2 — CERTIFICATE
% ════════════════════════════════════════════════════════════
\chapter*{}
\thispagestyle{fancy}
\vspace*{-1.5cm}
\begin{center}
    {\normalsize PCET's \& NMVPM's}\\[0.2cm]
    {\large \textbf{NUTAN MAHARASHTRA INSTITUTE OF ENGINEERING \& TECHNOLOGY}}\\
    {\normalsize TALEGAON DABHADE, PUNE -- 410507}\\[0.5cm]

    % ── Replace 'nmiet_logo.png' with your actual logo file ──
    \includegraphics[width=0.22\textwidth]{nmiet_logo.png}\\[0.5cm]

    {\Huge \textit{Certificate}}\\[0.8cm]
\end{center}

\noindent This is to certify that, \textbf{Aditya Shailesh Bharate (Roll No. 01)} have successfully 
completed the Internship entitled and ``\underline{AWS AI-POWERED CLOUD ENGINEER}'' 
under my supervision, in the partial fulfilment of Bachelor of Engineering in Computer 
Engineering of Savitribai Phule Pune University (SPPU).

\vspace{2cm}
\noindent Date: 05/03/2025 \hfill Place: Pune

\vspace{3cm}
\noindent
\begin{tabular}{p{4cm} p{4.5cm} p{4.5cm}}
    \textbf{Internship Guide} & \textbf{Internship Coordinator} & \textbf{HOD Computer Engineering}
\end{tabular}


% ════════════════════════════════════════════════════════════
%  PAGE 3 — ACKNOWLEDGEMENT
% ════════════════════════════════════════════════════════════
\chapter*{ACKNOWLEDGEMENT}
\addcontentsline{toc}{chapter}{Acknowledgement}
\thispagestyle{fancy}

I would like to express our deepest gratitude and appreciation to all those who have 
contributed to the successful completion of this Internship. It is with great pleasure that we 
acknowledge their invaluable support, guidance, and encouragement throughout this journey.

I thank our Internship guide \textbf{Prof. Ashwini Utture}, for her exceptional support, 
guidance, and patience throughout this journey. Her expertise and mentorship have been 
invaluable in shaping our understanding of the internship and overcoming challenges.

I would like to express our gratitude to our Internship coordinator \textbf{Prof. Nilesh Kamble} 
for their valuable suggestions, encouragement, and motivation that have contributed to the 
successful completion of this internship. We extend our appreciation to all the faculty members 
for their continuous support and encouragement during the course of this Internship.

I express our heartily gratitude towards \textbf{Dr. Rahul Patil}, Head of Department of 
Computer Engineering for guiding us to understand the work conceptually and also for 
providing necessary information and required resources with his constant encouragement to 
complete this internship work.

With deep sense of gratitude, I thank our principal \textbf{Dr. Pramod Patil} and Management 
of the NMIET for providing all necessary facilities and their constant encouragement and 
support. The encouragement and academic environment provided by our institution were 
crucial in fostering our creativity and enabling us to explore new ideas.

My sincere thanks to all the teaching and non-teaching staff members of the Computer 
Engineering Department for their assistance, cooperation, and provision of the required 
resources throughout this Internship. I am ending this acknowledgement with deep 
indebtedness to my friends who have helped us.

I sincerely acknowledge and appreciate the contributions of all those who directly or 
indirectly supported us in the completion of this Internship.


% ════════════════════════════════════════════════════════════
%  PAGE 4 — TABLE OF CONTENTS
% ════════════════════════════════════════════════════════════
\tableofcontents
\thispagestyle{fancy}


% ════════════════════════════════════════════════════════════
%  CHAPTER I — ORGANIZATIONAL STRUCTURE OF INDUSTRY
% ════════════════════════════════════════════════════════════
\chapter{ORGANIZATIONAL STRUCTURE OF INDUSTRY}

\section{Organizational Structure:}

% If you have an org chart image, insert it here:
% \begin{figure}[h!]
%     \centering
%     \includegraphics[width=0.7\textwidth]{org_chart.png}
%     \caption{Organizational Structure of Eduskills Foundation}
% \end{figure}

\section{Address of the Company:}

\noindent \textbf{Corp. Office:} B-7, Awfis - Allied House, Near AICTE, Nelson Mandela Marg, 
Vasant Kunj, New Delhi -- 110070

\noindent \textbf{Head Office:} \#806, DLF Cyber City, Technology Corridor, 
Bhubaneswar - 751024, Odisha.

\noindent \textbf{Name:} Eduskills Foundation

\noindent \textbf{Address (Head Office):} \#806, DLF Cyber City, Technology Corridor, 
Bhubaneswar -- 751024, Odisha

\noindent \textbf{CIN:} [Enter CIN Number]

\section{Company Vision:}

Transforming the vision of skilled India and education for to benefit the education 
ecosystem by providing 360 degree holistic solutions to all the stakeholders.

\section{Company Mission:}

To positively impact 1 million beneficiaries by 2026.

\medskip
\noindent \textbf{How:}

By comprehensive identification of skills gaps in the students and mapping them with 
latest and world's best technical skills.


% ════════════════════════════════════════════════════════════
%  CHAPTER II — TECHNOLOGY LEARNED IN COMPANY
% ════════════════════════════════════════════════════════════
\chapter{TECHNOLOGY LEARNED IN COMPANY}

During the internship conducted through the Amazon Web Services Educate platform, 
several modern technologies related to cloud computing and artificial intelligence 
were studied. The training focused on understanding how cloud platforms help 
organizations store data, run applications, and manage infrastructure efficiently without 
relying on physical hardware. This approach improves scalability, cost efficiency, and 
system reliability, making cloud computing a key technology used by modern 
businesses and technology companies.

The internship also introduced practical knowledge of different AWS services used for 
computing, storage, networking, databases, serverless applications, and artificial 
intelligence. Through structured learning modules and assessments, a strong foundation 
was developed in cloud infrastructure, serverless architecture, machine learning 
fundamentals, and generative AI technologies, helping to understand how intelligent 
applications can be built using cloud platforms.

\medskip
\noindent \textbf{Technologies Learned}

\begin{itemize}[leftmargin=1.5cm]
    \item Cloud Computing Fundamentals
    \item Amazon S3 -- Cloud Storage Service
    \item Amazon EC2 -- Virtual Cloud Computing Instances
    \item Virtual Private Cloud (VPC) -- Cloud Networking
    \item Amazon RDS -- Managed Relational Database Service
    \item AWS Lambda -- Serverless Computing
    \item Amazon SageMaker -- Machine Learning Platform
    \item Amazon Bedrock -- Generative AI Service
    \item Cloud Security and Monitoring Concepts
    \item AI-Powered Cloud Application Development
\end{itemize}


% ════════════════════════════════════════════════════════════
%  CHAPTER III — DESCRIPTION OF INTERNSHIP
% ════════════════════════════════════════════════════════════
\chapter{DESCRIPTION OF INTERNSHIP}

During my internship, I gained foundational knowledge and practical exposure to cloud 
computing and artificial intelligence technologies through the Amazon Web Services 
Educate learning platform. The training consisted of structured modules designed to 
introduce key cloud concepts and widely used AWS services. Through these modules, I 
explored how cloud platforms enable organizations to store data, deploy applications, 
and manage infrastructure efficiently without relying on traditional on-premise systems.

The internship was organized into multiple learning modules covering topics such as 
cloud fundamentals, storage services, compute services, networking, databases, cloud 
security, serverless computing, machine learning, and generative AI. Each module 
included instructional materials, demonstrations, and assessments that helped me 
understand the working principles of different cloud services and their role in building 
modern cloud-based applications.

During the training, I explored several important AWS services such as Amazon S3 for 
cloud storage, Amazon EC2 for virtual computing resources, and AWS Lambda for 
serverless computing. Learning about these technologies helped me understand how 
developers can build applications that automatically scale based on demand without 
managing physical infrastructure. I also studied database management using Amazon 
RDS, which simplifies database deployment and maintenance in cloud environments.

Additionally, the internship introduced emerging technologies such as machine 
learning and generative artificial intelligence using tools like Amazon SageMaker 
and Amazon Bedrock. These modules provided insight into how AI models can be 
trained, deployed, and integrated with cloud services to develop intelligent applications. 
Overall, the internship broadened my understanding of modern cloud ecosystems and 
gave me practical exposure to technologies that are widely used in today's technology 
industry.


% ════════════════════════════════════════════════════════════
%  CHAPTER IV — DESCRIPTION OF TASKS
% ════════════════════════════════════════════════════════════
\chapter{DESCRIPTION OF TASKS}

During the internship, I completed a structured set of learning modules through the 
Amazon Web Services Educate platform as part of the AI Powered Cloud Engineer 
training program. Each module focused on a different aspect of cloud computing and 
artificial intelligence technologies that are widely used in modern IT infrastructure.

The initial module introduced the fundamentals of cloud computing, including the 
concept of cloud infrastructure, benefits of cloud adoption, and different service models 
used by cloud providers. The next modules focused on cloud storage and computing 
services, where I studied how Amazon S3 is used for storing and managing large 
amounts of data and how Amazon EC2 provides scalable virtual computing resources.

Further modules covered cloud networking and databases, where I explored the concept 
of Virtual Private Cloud (VPC) and learned how organizations manage relational 
databases using Amazon RDS. Additional modules introduced cloud operations and 
security, focusing on monitoring systems, access management, and protection of cloud 
resources.

The internship also included advanced topics such as serverless computing using AWS 
Lambda, machine learning foundations using Amazon SageMaker, and generative 
artificial intelligence using Amazon Bedrock. These modules helped me understand how 
cloud platforms can support intelligent and scalable applications.
\newpage
\section{Introduction of Projects:}

The main objective of this internship project was to gain practical knowledge of cloud 
computing technologies and understand how cloud platforms support modern software 
systems. With the rapid growth of digital services, cloud computing has become an 
essential technology used by organizations to store data, deploy applications, and 
manage computing infrastructure efficiently.

This project focused on exploring the services and architecture provided by Amazon 
Web Services, one of the leading cloud computing platforms in the world. AWS 
provides a wide range of services that allow businesses to build scalable, secure, and 
cost-effective applications without the need to maintain physical hardware.

During the internship, I studied various AWS services that are used for different 
purposes in cloud environments. These included cloud storage services such as 
Amazon S3, which enables secure storage and management of large volumes of data, 
and computing services such as Amazon EC2, which allow users to run virtual servers 
in the cloud. I also explored networking technologies through Virtual Private Cloud 
(VPC), database management through Amazon RDS, and serverless computing using 
AWS Lambda.

In addition to cloud infrastructure technologies, the project also introduced artificial 
intelligence tools such as Amazon SageMaker for machine learning and Amazon 
Bedrock for generative AI applications. These technologies demonstrate how cloud 
computing platforms are evolving to support intelligent data-driven systems.

I built an event-driven serverless system on AWS to optimize cloud costs. 
The problem was that idle EC2 instances often continue running and increase expenses. 
To solve this, I used EventBridge to trigger a Lambda function on a schedule. The 
Lambda function checks EC2 instances and uses CloudWatch metrics to identify low 
CPU usage. If an instance is idle, it automatically stops it.

I also implemented tag-based filtering to ensure only non-production resources are 
affected and added logging for monitoring. One challenge was defining what counts as 
`idle', which I solved by setting CPU thresholds. This solution helps reduce unnecessary 
cloud spending and demonstrates a scalable, serverless approach to automation.
\newpage
\section{Motivation and Relational Study:}

The increasing use of digital platforms and online services has created a strong demand 
for scalable computing infrastructure. Cloud computing has emerged as a powerful 
solution that allows organizations to access computing resources such as storage, 
processing power, and databases through the internet without maintaining their own 
hardware systems.

The motivation behind undertaking this internship was to gain practical exposure to 
cloud technologies and understand how modern IT systems are designed using cloud 
platforms. Learning about cloud computing provides valuable insight into how 
businesses manage data, deploy applications, and handle large-scale workloads 
efficiently.

Another important motivation was to explore the integration of artificial intelligence 
with cloud computing. AI technologies such as machine learning and generative AI 
require large amounts of data and computing resources, which are efficiently supported 
by cloud platforms. Understanding how these technologies interact helps in developing 
modern intelligent systems.

This internship also helped bridge the gap between theoretical knowledge and real-world 
technology applications. Concepts learned in areas such as networking, databases, and 
software systems could be related to practical cloud services provided by AWS. By 
studying these technologies together, I developed a clearer understanding of how cloud 
infrastructure, data processing systems, and AI models work collectively to support 
modern digital applications.
\newpage
\section{Methodology Details:}

The methodology followed during the internship was based on a structured learning 
approach through the Amazon Web Services Educate training modules. Each module 
focused on a specific cloud technology and included learning materials and assessments 
to evaluate understanding. This step-by-step approach helped build a clear 
understanding of how different cloud services operate and how they are used in 
real-world applications.

The process began with studying the fundamentals of cloud computing, including cloud 
architecture, service models, and the benefits of cloud adoption such as scalability, 
reliability, and cost efficiency. I also learned about different cloud service models like 
Infrastructure as a Service (IaaS), Platform as a Service (PaaS), and Software as a 
Service (SaaS).

The next modules focused on cloud storage using Amazon S3, where I learned how 
large volumes of data can be securely stored and managed in the cloud. This was 
followed by studying cloud computing resources through Amazon EC2, which 
explained how virtual machines are deployed and scaled based on application 
requirements.

Further modules introduced cloud networking through Virtual Private Cloud (VPC) 
and database management using Amazon RDS, which demonstrated how secure 
networks and relational databases are managed in cloud environments.

Another important module introduced serverless computing using AWS Lambda. 
Serverless computing allows developers to run application code without managing 
underlying servers. The module explained how event-driven architecture works and how 
applications can automatically scale based on incoming requests. This concept 
highlighted the efficiency and flexibility of modern cloud-based application development.

The final modules focused on machine learning and artificial intelligence technologies. 
Using Amazon SageMaker, I studied the basic workflow of machine learning including 
data preparation, model training, and deployment. Additionally, the internship 
introduced generative artificial intelligence using Amazon Bedrock, which provides 
access to advanced AI models capable of generating text and other forms of content.


% ════════════════════════════════════════════════════════════
%  CHAPTER V — LEARNING OUTCOMES
% ════════════════════════════════════════════════════════════
\chapter{LEARNING OUTCOMES}

Through this internship, I gained a clear understanding of the fundamental concepts of 
cloud computing and artificial intelligence technologies. The training helped me 
understand how cloud platforms provide scalable infrastructure for storing data, running 
applications, and managing computing resources efficiently. I also learned about 
different cloud service models and how organizations use cloud services to improve 
performance, reliability, and cost efficiency.

During the internship, I developed knowledge of several important cloud services 
provided by Amazon Web Services. I learned how Amazon S3 is used for secure and 
scalable cloud storage, and how Amazon EC2 provides virtual computing resources for 
deploying applications. I also understood the importance of cloud networking, database 
management using Amazon RDS, and serverless computing through AWS Lambda.

In addition, the internship introduced me to emerging technologies such as machine 
learning using Amazon SageMaker and generative artificial intelligence using Amazon 
Bedrock. Overall, the internship helped me develop a strong foundation in cloud 
technologies and provided insight into how AI and cloud computing can be combined 
to build modern intelligent applications.


% ════════════════════════════════════════════════════════════
%  CHAPTER VI — CONCLUSION
% ════════════════════════════════════════════════════════════
\chapter{CONCLUSION}

The internship provided valuable exposure to the fundamentals of cloud computing and 
artificial intelligence technologies through the learning modules on the Amazon Web 
Services Educate platform. Throughout the internship, I studied various cloud services 
and gained an understanding of how modern cloud infrastructure supports the 
development and deployment of scalable applications. The training modules helped me 
understand the architecture and working principles of different cloud technologies.

During the internship, I explored important AWS services such as Amazon S3 for cloud 
storage, Amazon EC2 for computing resources, Amazon RDS for database management, 
and AWS Lambda for serverless computing. I also gained introductory knowledge of 
machine learning and generative AI using Amazon SageMaker and Amazon Bedrock. 
These technologies demonstrated how cloud platforms enable efficient data processing, 
application deployment, and intelligent system development.

Overall, the internship strengthened my understanding of cloud-based technologies and 
provided insight into how cloud computing and artificial intelligence are used in modern 
IT environments. The knowledge gained during this training will help in further 
developing technical skills and exploring advanced technologies in the field of cloud 
computing and AI.


% ════════════════════════════════════════════════════════════
%  CHAPTER VII — REFERENCES
% ════════════════════════════════════════════════════════════
\chapter*{REFERENCES}
\addcontentsline{toc}{chapter}{References}
\thispagestyle{fancy}

\begin{center}
    \rule{\linewidth}{0.4pt}
\end{center}

\begin{enumerate}[leftmargin=1.5cm, label={[\arabic*]}]

    \item Amazon Web Services, \textit{AWS Lambda -- Serverless Computing Documentation}. [Online]. Available: \url{https://docs.aws.amazon.com/lambda/}

    \item Amazon Web Services, \textit{Amazon SageMaker -- Machine Learning Documentation}. [Online]. Available: \url{https://docs.aws.amazon.com/sagemaker/}

    \item Amazon Web Services, \textit{Amazon Bedrock -- Generative AI Documentation}. [Online]. Available: \url{https://docs.aws.amazon.com/bedrock/}

    \item Amazon Web Services, \textit{Amazon EventBridge Documentation}. [Online]. Available: \url{https://docs.aws.amazon.com/eventbridge/}

    \item Amazon Web Services, \textit{Amazon CloudWatch -- Monitoring and Observability}. [Online]. Available: \url{https://docs.aws.amazon.com/cloudwatch/}

    \item E. Jonas et al., ``Cloud programming simplified: A Berkeley view on serverless computing,'' \textit{arXiv preprint}, arXiv:1902.03383, Feb. 2019. [Online]. Available: \url{https://arxiv.org/abs/1902.03383}

    \item T. Brown et al., ``Language models are few-shot learners,'' in \textit{Proc. Advances in Neural Information Processing Systems (NeurIPS)}, vol. 33, 2020, pp. 1877--1901. [Online]. Available: \url{https://arxiv.org/abs/2005.14165}

    \item A. Vaswani et al., ``Attention is all you need,'' in \textit{Proc. NeurIPS}, vol. 30, 2017, pp. 5998--6008. [Online]. Available: \url{https://arxiv.org/abs/1706.03762}

    \item W. Zhao et al., ``A survey of large language models,'' \textit{arXiv preprint}, arXiv:2303.18223, 2023. [Online]. Available: \url{https://arxiv.org/abs/2303.18223}

    \item K. Ren, C. Wang, and Q. Wang, ``Security challenges for the public cloud,'' \textit{IEEE Internet Computing}, vol. 16, no. 1, pp. 69--73, Jan.--Feb. 2012, doi: 10.1109/MIC.2012.14.

    \item S. Gill et al., ``Transformative effects of ChatGPT on modern education: Emerging era of AI chatbots,'' \textit{Internet of Things and Cyber-Physical Systems}, vol. 4, pp. 19--23, 2024, doi: 10.1016/j.iotcps.2023.06.002.

\end{enumerate}


% ════════════════════════════════════════════════════════════
%  ANNEXURE A — NOC LETTER
% ════════════════════════════════════════════════════════════
% ════════════════════════════════════════════════════════════
%  ANNEXURE A — NOC LETTER
% ════════════════════════════════════════════════════════════
\appendix
\chapter*{Annexure A}
\addcontentsline{toc}{chapter}{Annexure A -- No Objection Certificate}
\thispagestyle{fancy}

\begin{center}
    {\small Nutan Maharashtra Vidya Prasarak Mandal's (NMVPM's)}\\[0.2cm]
    {\large \textbf{NUTAN MAHARASHTRA INSTITUTE OF ENGINEERING AND TECHNOLOGY (NMIET)}}\\
    {\small Under Administrative Support - Pimpri Chinchwad Education Trust (PCET)}\\[0.2cm]
    {\small Approved by AICTE \quad Accredited by NAAC \& NBA \quad Affiliated to SPPU}\\[0.2cm]
    {\large \textbf{DEPARTMENT OF COMPUTER ENGINEERING}}\\
\end{center}

\vspace{0.3cm}
\noindent \textbf{Ref No.:} \hfill \textbf{Date:}

\vspace{0.3cm}
\noindent To\\
\textbf{The HR Manager,}

\vspace{0.3cm}
\noindent \textbf{Subject: No Objection Certificate for Internship}

\vspace{0.2cm}
\noindent Respected Sir/Madam,

\vspace{0.2cm}
\noindent This is to certify that \underline{\hspace{3cm}} (PRN: \underline{\hspace{2cm}}), a bonafide student of \textbf{Third Year (Branch: Computer Engineering)}, studying in \textbf{Semester VI} at \textbf{Nutan Maharashtra Institute of Engineering and Technology, Pune, Department of Computer Engineering}, affiliated to \textbf{Savitribai Phule Pune University}, Pune, is permitted to undergo \textbf{Internship} at your esteemed organization.

\vspace{0.2cm}
\noindent The institute has \textbf{no objection} to the student undergoing internship from \textbf{Start Date:} \underline{\hspace{2.5cm}} to \textbf{End Date:} \underline{\hspace{2cm}} mode as per the company guidelines. This internship is a part of the academic curriculum of SPPU, Pune. The student will abide by all the rules, regulations, and confidentiality policies of your organization during the internship period. We shall be thankful for your kind cooperation and support.

\vspace{0.8cm}
\noindent
\begin{tabular}{p{3.5cm} p{4.5cm} p{4.5cm}}
     & Prof. Nilesh Kamble & Dr. Rahul Patil \\[0.3cm]
    \textbf{Internship Guide} & \textbf{Internship Coordinator} & \textbf{HOD}\\
     & & \textbf{Computer Engineering}
\end{tabular}


% ════════════════════════════════════════════════════════════
%  ANNEXURE B — COMPLETION CERTIFICATE
% ════════════════════════════════════════════════════════════
\chapter*{Annexure B}
\addcontentsline{toc}{chapter}{Annexure B -- Completion Certificate}
\thispagestyle{fancy}

\begin{center}
    {\Large \textbf{COMPLETION CERTIFICATE}}
\end{center}

\vspace{0.3cm}

\begin{center}
\includegraphics[width=0.85\textwidth, height=0.65\textheight, keepaspectratio]{completion_certificate.png}
\end{center}

\vspace{0.3cm}


% ════════════════════════════════════════════════════════════
%  LOG BOOK
% ════════════════════════════════════════════════════════════
\chapter*{Log Book}
\addcontentsline{toc}{chapter}{Log Book}
\thispagestyle{fancy}

\begin{center}
\small
\begin{tabular}{|p{1.2cm}|p{9cm}|p{2.8cm}|}
    \hline
    \textbf{Week} & \textbf{Work Done} & \textbf{Coordinator Sign} \\
    \hline
    Week 1 & Studied Cloud Computing Fundamentals -- cloud architecture, service models (IaaS, PaaS, SaaS) and benefits of cloud adoption via AWS Educate. & \\[0.4cm]
    \hline
    Week 2 & Explored Amazon S3 -- cloud storage concepts, bucket creation, object storage and access policies. & \\[0.4cm]
    \hline
    Week 3 & Studied Amazon EC2 -- virtual computing instances, instance types, launch configurations and application deployment. & \\[0.4cm]
    \hline
    Week 4 & Learned VPC and Amazon RDS -- cloud networking, subnets, security groups and managed relational database deployment. & \\[0.4cm]
    \hline
    Week 5 & Studied Cloud Security and Monitoring -- IAM roles, access management, CloudWatch metrics, logging and protection of cloud resources. & \\[0.4cm]
    \hline
    Week 6 & Explored AWS Lambda and Serverless Computing -- event-driven architecture, function deployment and automatic scaling concepts. & \\[0.4cm]
    \hline
    Week 7 & Studied Amazon SageMaker -- machine learning workflow including data preparation, model training, evaluation and cloud deployment. & \\[0.4cm]
    \hline
    Week 8 & Explored Amazon Bedrock and Generative AI -- large language models, generative AI capabilities and completed all module assessments. & \\[0.4cm]
    \hline
    Week 9 & Built an event-driven serverless system on AWS using EventBridge and Lambda to auto-stop idle EC2 instances based on CloudWatch CPU metrics. Implemented tag-based filtering for non-production resources and CPU threshold-based idle detection to reduce cloud costs. & \\[0.4cm]
    \hline
\end{tabular}
\end{center}

\end{document}